\subsection{Robustness}

We now review the notion of robustness, building upon the formalism presented by Gurkan et al.~\cite{GurkanJMMST21}.
In the literature, robustness today refers to the property that so long as the number of corrupted and honest participants meet some pre-defined threshold, then the protocol will complete successfully.
The definition by Gurkan et al.\ deems a DKG to be robust if
honest parties' shares combine to output a secret key that is valid with respect to a public key.
However, the notion by Gurkan et al.\ does not consider the consistency of players' ending states.

We present several clarifications to the notion of robustness using a game-based definition in Figure~\ref{fg:robustness}.
First, our definition is generalizable beyond discrete-logarithm based constructions.
Second, our definition encodes two different forms of robustness.
\emph{Strong robustness} defines the existing notion of robustness in the literature.
We additionally introduce the notion of \emph{weak robustness},
which allows the protocol to non-fatally abort,
so long as all honest participants maintain the same status.
In particular, participants can exit either with an $\abort$ status or a $\fail$ status,
where $\fail$ indicates a critical failure.
A critical failure depends on the context of the construction;
for the generic construction we present in Section~\ref{generic-construction},
a critical failure occurs when the adversary can cause the protocol to abort in
the final round of the protocol, which in turn has consequences for the proof
of security.

Prior notions in the literature considered only strong robustness;
however, introducing a weaker variant allows us to define a more efficient construction that fulfills our other notions of secrecy and indistinguishability.

\begin{remark}[Number of Honest Players]
  The robustness experiment in Figure~\ref{fg:robustness} assumes an honest majority of players.
  If an adversary were allowed to control the majority of  participants,
  it could split the view of honest participants,
  unless some additional structure is assumed, such as a public bulletin board.
\end{remark}

\begin{figure}[p]
  \begin{pchstack}[boxed,center,space=1em]
  \begin{pcvstack}
    \procedure[]{\protect\fbox{$\StrongRobustnessExp_{\dkg,\adv}(\lambda)$} \protect\dbox{$\WeakRobustnessExp_{\dkg,\adv}(\lambda)$} }{
    \roundcounter \gets 0 \text{\algcom Ensure rounds are queried in order}\\
    \invecp_{11} \gets \emptyset, \ldots, \invecp_{1n} \gets \emptyset \text{\algcom Simulate peer-to-peer channels} \\
    (n, t, \corrupt, \advstate) \gets \adv(1^\lambda);\
    \pcreturn 0\ \pcif n < t \vee \corrupt \not\subset \set{n} \vee \lvert \corrupt \rvert \ge t \\
    \honest \gets \set{n} \setminus \corrupt \\
    \pcreturn 0\ \pcif \lvert \honest \rvert < t \\
    \pcfor i \in \honest \pcdo \\
    \pcind (\party_i, \outvecp_{0i}, \bmsg_i) \gets \dkgkeygen_0(\lambda,n, t, i) \\
    \pcind \invecp_{1k}[i] \gets \outvecp_{0i}[k],\ \forall\ k \in \set{n}\\
    \text{\algcom Set party $i$'s peer-to-peer messages for all other parties} \\
    \adv^{{\oraclekg}, \oraclefinalize}(\advstate, \{ \invecp_{1i} \}_{i \in \corrupt}, \{ \bmsg_i \}_{i \in \honest}) \\
    \pcreturn 0\ \pcif \roundcounter \neq \numrounds \text{\algcom Prevents trivial win by forcing coordination} \\
    \fbox{$ \pcif \exists i \in \honest \ : \ \status_i \neq \accept $} \\
    %\fbox{$\pcif \honest \neq \{i : \status_i = \accept \}_{i \in \honest}$} \\
    \fbox{$\pcind \pcreturn 1$} \\
    \fbox{$ \pcif \exists i, j \in \honest \ : \ \status_i \neq \status_j $} \\
    \fbox{\pcind \pcreturn 1} \\
    %\dbox{$\pcfor i,j \in \honest \pcdo$ } \\
    %\dbox{$\pcind \pcreturn 1\ \pcif \status_i \neq \status_j $} \\
    \dbox{\algcom $\adv$ wins if any honest party ends in a disjoint state} \\
    \dbox{$ \pcif \forall i \in \honest\ :\ \status_i = \abort $} \\
    %\dbox{$\pcif \honest = \{ i : \status_i = \abort \}_{i \in \honest}$} \\
    \dbox{$\pcind \pcreturn 0 $} \\
    \dbox{\algcom $\adv$ loses if all honest parties end in a consistent aborted state } \\
    \pcfor i,j \in \honest \pcdo  \\
    \pcind \pcreturn 1\ \pcif \PK_i \neq \PK_j  \\
    \pcind \pcreturn 1\ \pcif \qual_i \neq \qual_j  \\
    \pcind \pcreturn 1\ \pcif \auxinfo_i \neq \auxinfo_j  \\
    \PK = \PK_i, i \in \honest \text{\algcom All honest parties at this point have the same view of $\pk$} \\
    \SK \gets \dkg.\recover(\PK, \{(i,\SK_i) \}_{i \in \honest}) \\
    \pcreturn 1\ \pcif \SK = \fail \\
    %\text{\algcom $\adv$ wins if honest parties' shares do not recover a valid secret} \\
    \pcreturn 0
  }
  \end{pcvstack}
  \end{pchstack}

  \caption[Robustness experiment for a DKG.]{Robustness experiment for a DKG.
Weak robustness differences are highlighted in a dashed box;
strong robustness differences are highlighted in a solid box.
}
%\douglas{Line 24 checks that $\sk \ne \fail$.  Should we also be checking that $\sk$ is the (a) secret key corresponding to $\pk$?  Is robustness violated if the adversary the parties can be tricked into recovering a different $\sk$?  Sorry if we've discussed this already and I forgot.}
% \chelsea{The recover algorithm checks key validity internally}
\label{fg:robustness}
\end{figure}

\begin{figure}[p]
  \begin{framed}
  %\begin{pchstack}[boxed,center,space=1em]
  \begin{pchstack}
  {
    %\chelsea{TODO: this needs to be script size for the proceedings version}
  }
	\pseudocode{
    \procedure[]{$\oraclekg(\{\{ \pmsg_{i,j}\}_{j \in \honest}, \bmsg_{ri} \}_{i \in \corrupt})
    $}{
      \text{\algcom Executes round  } \\
      \text{\algcom $r \in \set{\numrounds-1}$ } \\
    \pcreturn \bot\ \pcif r \neq \roundcounter+1 \\
    \text{\algcom Ensure all prior rounds } \\
    \text{\algcom have been queried,} \\
    \text{\algcom and that this round has} \\
    \text{\algcom not yet been queried.} \\
    \roundcounter \gets r \\
    \pcfor i \in \corrupt, j \in \honest \pcdo \\
      \pcind \invecp_{rj}[i]  \gets \pmsg_i[j] \\
    \pcfor j \in \honest \pcdo \\
      \pcind \mathsf{in} \gets (\party_j, \invecp_{rj}, \{\bmsg_{ri} \}_{i \in \set{n}}) \\
      \pcind \fbox{$\mathsf{out} \gets \dkgkeygen_r(\mathsf{in})$} \\
      \pcind \dbox{$\mathsf{out} \gets \simkeygen_r(\mathsf{in})$} \\
      \pcind (\party_j, \outvecp_{r+1,j}, \bmsg_{r+1,j}) \gets \mathsf{out} \\
      \pcind \pcfor k \in \set{n} \pcdo \\
      \pcind \pcind \invecp_{r+1,k} \gets \outvecp_{r+1,j}[k] \\
    \pcreturn (\{  \invecp_{r+1,i} \}_{i \in \corrupt}, \{ \bmsg_j \}_{j \in \honest} )
  }
  \pchspace
  \begin{pcvstack}
  \procedure[]{$\oraclefinalize(\{\{ \pmsg_{i,j}\}_{j \in \honest}, \bmsg_{ri} \}_{i \in \corrupt})$}{
    \pcreturn \bot\ \pcif \roundcounter \neq \numrounds - 1 \\
    \roundcounter \gets \roundcounter + 1 \\
    \pcfor i \in \corrupt, j \in \honest \pcdo \\
      \pcind \invecp_{rj}[i]  \gets \pmsg_i[j] \\
    \pcfor j \in \honest \pcdo \\
      \pcind \mathsf{in} \gets (\party_j, \invecp_{rj}, \{\bmsg_{ri} \}_{i \in \set{n}}) \\
      \pcind \fbox{$\mathsf{out} \gets \finalize(\mathsf{in}) $}\\
      \pcind \dbox{$\mathsf{out} \gets \simfin(\mathsf{in}) $}\\
      \pcind (\PK_j, \qual_j, \auxinfo_j), (\status_j, \SK_j) \gets \mathsf{out}
  }
  \end{pcvstack}
  }
  \end{pchstack}
  \end{framed}
\caption[Common oracles for DKG security experiemnts.]{Common oracles used within security experiments for a distributed key generation scheme.
Boxed algorithms are used by the robustness and indistinguishability games;
dashed boxed algorithms are used by the zero-knowledge game.
}
\label{fg:keygen-o}
\end{figure}




\paragraph{Common Oracles.}
We define common oracles that are used within the respective security
experiments for a DKG in in Figure~\ref{fg:keygen-o}.
Boxed algorithms are used by the robustness and indistinguishability games;
dashed boxed algorithms are used by the zero-knowledge game.

\subsubsection{Strong Robustness}
We show the strong robustness attack experiment in Figure~\ref{fg:robustness},
where its differences with weak robustness are outlined in a box.


The advantage of an adversary $\adv$ against $\dkg$ in the strong robustness experiment is
  \[
  \advantstrrobust{\dkg,\adv}(\lambda) = \Pr[\StrongRobustnessExp_{\dkg,\adv}(\lambda) = 1 ]
\]

\begin{definition}
A DKG $\dkg$ is \defn{strongly robust} if for all probabilistic polynomial-time adversaries $\adv$,
the function $\advantstrrobust{\dkg,\adv}(\lambda)$  is negligible.
\end{definition}


The strong robustness attack game allows $\adv$ to participate in key generation and control up to $t-1$ players,
the set of which we denote by $\corrupt$.
We denote the set of honest players as $\honest$.
The robustness attack game has oracles $\oraclekg, \currentround \in
\{ 0, \ldots, \numrounds-1 \}  $
and $\oraclefinalize$ in common with other games,
so we show these oracles in Figure~\ref{fg:keygen-o}.
% This requirement doesn't seem necessary in the definition
% and require that the number of honest players be at least $t$.
The adversary participates in the protocol by exchanging peer-to-peer and
broadcast messages with honest players in each round $\currentround$ via the oracle
$\oraclekg$.
The experiment is responsible for exchanging messages between parties and enforcing assumptions of the underlying network channels.
Additionally, the experiment ensures that rounds must be queried in order, and that no round is repeated.

$\adv$ wins the strong robustness game under three conditions.
First, if it causes any honest player to end in a non-accepting state.
Second, if it causes any honest player's view of $\pk, \qual,$ or $\auxinfo$ to
be different from any other honest player's view.
Finally, $\adv$ wins if it can cause the resulting keypair to be invalid with
respect to the target key generation algorithm (recall that $\dkg.\recover$ tests for validity of the keypair, and outputs $\fail$ if $\sk$ is invalid).

\subsubsection{Weak Robustness}
We similarly present the weak robustness attack experiment in Figure~\ref{fg:robustness},
where its differences with strong robustness are outlined in a dashed box.
Intuitively, weak robustness differs from strong robustness in that it
allows honest parties to end in an aborted state $\abort$,
so long as \emph{all} honest parties end in the same state.
Otherwise, if all honest parties end in an accepting state,
it is identical to the notion of strong robustness.

The advantage of an adversary $\adv$ against the DKG $\dkg$ in the weak robustness experiment is
\[
  \advantwkrobust{\dkg,\adv}(\lambda) = \Pr[\WeakRobustnessExp_{\dkg,\adv}(\lambda) = 1 ]
\]

\begin{definition}
A DKG $\dkg$ is \defn{weakly robust} if for all probabilistic polynomial-time adversaries $\adv$,
the function $\advantwkrobust{\dkg,\adv}(\lambda)$  is negligible.
\end{definition}


$\adv$ wins the weak robustness game under three conditions.
First, $\adv$ wins if honest players do not all finish with the same status.
Otherwise, if all honest players have status $\accept$, then $\adv$ wins
if the honest players' view of $\pk$ or $\qual$ are not all the same or if
the resulting keypair is invalid with respect to the target key generation algorithm.



